\documentclass[11pt]{article}

\usepackage[margin=1in]{geometry}
\usepackage{amsmath,amssymb,amsthm,amsfonts}
\usepackage{mathtools}
\usepackage{bm}
\usepackage{hyperref}
\usepackage{enumitem}
\usepackage[T1]{fontenc}
\usepackage{lmodern}
\usepackage{microtype}
\usepackage{booktabs}
\usepackage{physics}
\usepackage{xcolor}
\usepackage{listings}

\hypersetup{
  colorlinks=true,
  linkcolor=blue!60!black,
  citecolor=blue!60!black,
  urlcolor=blue!60!black
}

\lstdefinestyle{pytorch}{
  language=Python,
  basicstyle=\ttfamily\small,
  keywordstyle=\color{blue!70!black}\bfseries,
  commentstyle=\color{green!40!black},
  stringstyle=\color{red!60!black},
  showstringspaces=false,
  frame=single,
  breaklines=true,
  tabsize=2
}

\title{Prime-Indexed Multiplicity Dynamics:\\
From Multiplicity Operator Calculus to MultiplicityCell Prime-Mixing Blocks}

\author{Citizen Gardens \\ The Foundation of Multiplicity}

\date{April 16, 2026}

\begin{document}

\maketitle

\begin{abstract}
This report synthesizes the developments discussed in the present thread on
Multiplicity Theory, focusing on the relationship between the Multiplicity
Operator Calculus (MOC) and a finite-dimensional neural surrogate,
\emph{MultiplicityCell}. The guiding idea is that reality---physical,
cognitive, computational---is governed by prime-indexed multiplicity dynamics
on a lawful Hilbert space, with a universal operator decomposition and
constitutional projectors enforcing lawfulness and ethics. Within this view,
MOC provides a detailed, prime-indexed, noncommutative operator calculus on
hypergraphs and cyclic time, while MultiplicityCell instantiates a compressed,
trainable prime-indexed recursion (with blocks \(A,B,E\)) suitable for
implementation in neural architectures. We give an executive summary, a
compact constitutional overview, a mapping from MOC to the global framework,
and a mathematically precise definition of the prime-mixing block \(A\) for
MultiplicityCell, including a concrete code sketch.
\end{abstract}
\newpage
\tableofcontents
\newpage
\section{Executive Summary}

The discussion across this thread clarified and refined a unifying view of
Multiplicity Theory:

\begin{itemize}[leftmargin=1.5em]
  \item There is a \emph{single constitutional backbone}: a lawful Hilbert
  space \(H_{\text{lawful}} \subset \ell^2(\mathbb{P}) \otimes L^2(\mathbb{R}) \otimes \mathbb{C}^d\)
  equipped with a universal operator
  \[
    U = A + B + E,
  \]
  together with a constitutional projector \(\Pi_{\text{CSL}}\) and an ethical
  projector \(P_E\), enforcing lawfulness and ethical viability of trajectories.

  \item \emph{Prime-indexed multiplicity} is the mother structure: primes
  index irreducible channels; composite structure arises via tensor products,
  factorization, and Chinese remainder decomposition. In this framework,
  prime-indexed operators generate the dynamics; composites are built via
  noncommutative words.

  \item The \emph{Multiplicity Operator Calculus} (MOC) \cite{MOC} provides an
  explicit, finite, discrete realization: a multiplicity space
  \(M = (X_n,H)\) with \(X_n\) signals on a cyclic time lattice and \(H\) a
  hypergraph of relations. Prime-indexed operator families
  \(\hat P_p\) act on both state and relation layers and are analyzed via a
  resonance functional \(R\in[0,1]\) combining time correlation, harmonic
  lock, and phase coherence. \cite{MOC}

  \item \emph{MultiplicityCell} is the neural (or numerical) surrogate of the
  universal recursion
  \[
    \psi_{t+1}
    =
    P_E\bigl(
      \Pi_{\text{CSL}}
      T_{\Lambda_m}(\psi_t, x_t)
    \bigr),
  \]
  where \(T_{\Lambda_m}\) decomposes into a prime-mixing block \(A\), a
  time-sieve block \(B\), and an internal block \(E\). The universal
  multiplicity constant \(\Lambda_m\) controls contraction and Absolute
  Contraction Energy (ACE) budgets.

  \item This report focuses on making the correspondence
  \(\text{MOC} \leftrightarrow (H_{\text{lawful}},U,\Pi_{\text{CSL}},P_E)\)
  explicit, and on defining the \emph{prime-mixing block} \(A\) within
  MultiplicityCell in a way that is faithful to MOC's structure while being
  implementable in neural architectures.
\end{itemize}

The upshot is that MOC can be viewed as a detailed local chart of the global
Multiplicity framework, and MultiplicityCell can be viewed as a compressed,
trainable instantiation of the same constitutional dynamics.

\section{Constitutional Overview of the Framework}

\subsection{Lawful Hilbert Space and Universal Operator}

We posit a lawful Hilbert space
\[
  H_{\text{lawful}} \subset \ell^2(\mathbb{P}) \otimes L^2(\mathbb{R}) \otimes \mathbb{C}^d,
\]
where:
\begin{itemize}[leftmargin=1.5em]
  \item \(\ell^2(\mathbb{P})\) indexes prime channels;
  \item \(L^2(\mathbb{R})\) (or a time/frequency surrogate) encodes temporal or
  spectral structure;
  \item \(\mathbb{C}^d\) carries internal degrees of freedom (e.g., relational,
  semantic, or gauge-like components).
\end{itemize}

The universal operator \(U\) is decomposed as
\[
  U = A + B + E,
\]
where, at an abstract level:
\begin{itemize}[leftmargin=1.5em]
  \item \(A\) is a prime-indexed mixing operator on \(\ell^2(\mathbb{P})\),
  potentially banded or otherwise constrained by the arithmetic structure of
  primes.
  \item \(B\) is a time-sieve or spectral operator acting primarily along the
  \(L^2(\mathbb{R})\) axis, akin to Fourier multipliers or filters over time
  and log-frequency.
  \item \(E\) acts on internal degrees of freedom, capturing recursion, local
  nonlinearity, or ethical/meta-constraints in the internal space.
\end{itemize}

The dynamics are governed by a recursion of the form
\begin{equation}
  \psi_{t+1}
  =
  P_E\bigl(
    \Pi_{\text{CSL}}\, T_{\Lambda_m}(\psi_t, x_t)
  \bigr),
  \label{eq:universal-recursion}
\end{equation}
where:
\begin{itemize}[leftmargin=1.5em]
  \item \(\Pi_{\text{CSL}}\) is a constitutional projector enforcing structural
  lawfulness (e.g., spectral, gauge, or tier-based constraints).
  \item \(P_E\) is an ethical projector enforcing viability and fairness
  constraints.
  \item \(T_{\Lambda_m}\) is a parametrized map incorporating the decomposition
  \(A+B+E\) and the universal multiplicity constant \(\Lambda_m\), which
  controls contraction and ensures the existence and uniqueness of lawful
  fixed points under appropriate conditions.
\end{itemize}

\subsection{Prime-Indexed Multiplicity}

Prime-indexed multiplicity enters in several ways:
\begin{itemize}[leftmargin=1.5em]
  \item Primes index channels in \(\ell^2(\mathbb{P})\).
  \item Operators are indexed by primes and their powers, reflecting the
  factorization structure of composite quantities.
  \item Chinese remainder theory (CRT) yields a tier decomposition of the
  time/frequency axis, with projectors onto prime-power levels.
  \item The decomposition of \(U\) respects prime grading: the \(A\)-block is
  prime-diagonal, while \(B\) and \(E\) respect prime-based constraints,
  possibly coupling primes in structured ways.
\end{itemize}

\section{Multiplicity Operator Calculus (MOC)}

\subsection{State Space and Operators}

MOC introduces a multiplicity space \(M = (X_n,H)\) where \cite{MOC}
\begin{itemize}[leftmargin=1.5em]
  \item \(T_n=\mathbb{Z}/n\mathbb{Z}\) is a cyclic time lattice;
  \item \(X_n = \{x:T_n \to \mathbb{R}^k\}\) is the space of signals or
  patterns on the lattice;
  \item \(H=(V,E,\iota)\) is a hypergraph encoding relational structure, with
  incidence map \(\iota:E\to\mathcal{P}(V)\) and state attachments on vertices
  and edges.
\end{itemize}

Prime-indexed operator families act on this space:
\begin{itemize}[leftmargin=1.5em]
  \item State operators on \(X_n\):
  \begin{itemize}
    \item Subdivision \(S_p\),
    \item Accent \(A_{p^r}\),
    \item Rotation \(R_{p^r}\),
    \item Permutation \(W_p\).
  \end{itemize}
  \item Relation operators on \(H\):
  \begin{itemize}
    \item Split, merge, fold, and relabel operators \(\hat Q_p\).
  \end{itemize}
\end{itemize}
Composite operator words \(W\) are ordered products of such prime-indexed
operators, acting noncommutatively on \((X_n,H)\) \cite{MOC}.

\subsection{Projectors, Tiers, and Resonance}

MOC introduces CRT-based projectors and spike gates. For each prime power
\(p^r\) dividing \(n\), one defines:
\begin{itemize}[leftmargin=1.5em]
  \item A level-averaging projector \(\Pi_{p^r}:X_n\to X_n\), averaging over
  positions separated by \(n/p^r\).
  \item A multiplicative gate \(\Delta_{p^r} x(t) = [p^r\,|\,t] x(t)\)
  selecting spikes at that tier.
\end{itemize}
These separate invariant content (via \(\Pi_{p^r}\)) from spike events (via
\(\Delta_{p^r}\)) at each prime-power level \cite{MOC}.

The resonance functional \(R\) measures the fit between a modeled pattern
\(x=W[x_0]\) and target data \(D\), with components
\begin{itemize}[leftmargin=1.5em]
  \item \(R_1\): time-domain correlation under optimal rotation;
  \item \(R_2\): harmonic lock using energy on comb indices corresponding to
  tiers;
  \item \(R_3\): phase coherence on tier-specific frequency subsets.
\end{itemize}
The aggregate score is
\[
  R = \lambda_1 R_1 + \lambda_2 R_2 + \lambda_3 R_3 \in [0,1],
\]
with \(\lambda_i \ge 0\), \(\sum_i \lambda_i=1\) \cite{MOC}.

\subsection{Relevance to the Global Framework}

Although MOC does not explicitly introduce \(H_{\text{lawful}}\), \(U\),
\(\Pi_{\text{CSL}}\), or \(P_E\), it exhibits structurally similar features:
\begin{itemize}[leftmargin=1.5em]
  \item A finite-dimensional surrogate of a prime-indexed space with time and
  relational degrees of freedom.
  \item Prime-indexed operators whose actions decompose along CRT tiers, akin
  to a prime block \(A\) and time-sieve block \(B\).
  \item Projectors \(\Pi_{p^r}\) and invariants (e.g., energy, fairness) that
  play a role analogous to constitutional and ethical projectors.
  \item A resonance functional \(R\) that certifies coherence, similar in
  spirit to ACE-based contractive certification schemes.
\end{itemize}

Thus, MOC may be viewed as a detailed ``chart'' of the global multiplicity
manifold, with its own natural operators and certification functionals.

\section{MultiplicityCell and the Prime-Mixing Block}

\subsection{MultiplicityCell as a Neural Surrogate}

MultiplicityCell is a finite-dimensional, trainable surrogate of the universal
recursion \eqref{eq:universal-recursion}. Its hidden state at time \(t\) can
be written as
\[
  h_t \in \mathbb{R}^{P \times d_{\text{int}}},
\]
where:
\begin{itemize}[leftmargin=1.5em]
  \item \(P\) is the number of prime channels retained (e.g., first \(P\)
  primes or a selected subset of primes),
  \item \(d_{\text{int}}\) is an internal dimension per prime channel.
\end{itemize}
We index primes as \(p_1,\dots,p_P\), with \(h_t[i,:]\) denoting the internal
vector for the \(i\)-th prime channel.

The core update step of the cell is implemented as
\[
  h_{t+1}
  =
  P_E\bigl(
    \Pi_{\text{CSL}}\, T_{\Lambda_m}(h_t, x_t)
  \bigr),
\]
where \(T_{\Lambda_m}\) is parametrized in terms of three blocks:
\begin{itemize}[leftmargin=1.5em]
  \item A \emph{prime-mixing} block \(A\) acting on prime-indexed channels;
  \item A \emph{time-sieve} block \(B\) acting along the temporal/frequency
  axis (usually realized as convolution or Fourier multipliers);
  \item An \emph{internal} block \(E\) acting on the internal dimension
  \(d_{\text{int}}\) and inputs \(x_t\).
\end{itemize}
The universal multiplicity constant \(\Lambda_m\) appears as a global
contraction parameter limiting the composite operator's Lipschitz constant.

\subsection{Design Requirements for the Prime-Mixing Block}

We now focus on the prime-mixing block \(A\). Its design should satisfy:
\begin{enumerate}[label=(\roman*), leftmargin=1.5em]
  \item \textbf{Prime-indexed structure:} The block should act on prime
  channels, respecting their indexing by \(\{p_i\}\).
  \item \textbf{Bandedness / local mixing:} The block should mix only nearby
  primes or primes linked by structural criteria (e.g., CRT tiers or semantic
  roles), mirroring MOC's structured cross-prime interactions \cite{MOC}.
  \item \textbf{Contraction:} The block should be contractive (or at least
  have controlled operator norm) to fit into ACE/Banach fixed-point arguments
  with a known contribution to \(\Lambda_m\).
  \item \textbf{Invariants and fairness:} The block should support or preserve
  invariants analogous to MOC's energy, tier energies, and fairness measures.
\end{enumerate}

\subsection{Neighborhood Structure on Prime Channels}

Let \(\{p_i\}_{i=1}^P\) be the distinct prime numbers indexing the channels.
We define a neighborhood system \(\mathcal{N}(i) \subseteq \{1,\dots,P\}\) for
each prime index \(i\), representing which primes may directly influence the
\(i\)-th channel. Choices include:
\begin{itemize}[leftmargin=1.5em]
  \item \emph{Local band in prime value:}
  \[
    \mathcal{N}(i) =
    \{\, j : |p_j - p_i| \le \Delta \,\}
  \]
  for some bandwidth parameter \(\Delta > 0\).
  \item \emph{CRT-tier based neighborhoods:} Use MOC's tier structure and
  projector hierarchy \(\Pi_{p^r}\) to connect primes that participate in the
  same or complementary tiers.
\end{itemize}

The neighborhood pattern encodes which prime channels may exchange information
directly through \(A\).

\subsection{Definition of the Prime-Mixing Block}

\begin{definition}[Prime-mixing block \(A\)]
Let \(h_t \in \mathbb{R}^{P \times d_{\text{int}}}\) be the hidden state at
time \(t\), indexed by primes \(p_1,\dots,p_P\). For each pair of indices
\((i,j)\) with \(1 \le i,j \le P\), let \(W_{ij} \in \mathbb{R}^{d_{\text{int}} \times d_{\text{int}}}\)
be a learnable matrix, and \(b_i \in \mathbb{R}^{d_{\text{int}}}\) a bias
vector. Fix neighborhood sets \(\mathcal{N}(i) \subseteq \{1,\dots,P\}\) for
all \(i\).

The \emph{prime-mixing block} \(A\) is the linear map
\[
  A : \mathbb{R}^{P \times d_{\text{int}}} \to \mathbb{R}^{P \times d_{\text{int}}}
\]
defined componentwise by
\begin{equation}
  (A h_t)_i
  =
  \sum_{j \in \mathcal{N}(i)}
    W_{ij} \, h_t[j,:] + b_i,
  \qquad i = 1,\dots,P.
  \label{eq:prime-mixing}
\end{equation}

We require that:
\begin{enumerate}[label=(\alph*), leftmargin=2em]
  \item \(W_{ij} = 0\) for all \(j \notin \mathcal{N}(i)\), i.e., the operator
  is banded with respect to the prime index neighborhoods.
  \item The operator norm satisfies
  \[
    \|A\|_{\mathrm{op}} \le \Lambda_A < 1,
  \]
  where \(\Lambda_A\) is a prescribed contraction parameter contributing to
  the global multiplicity constant \(\Lambda_m\).
\end{enumerate}
\end{definition}

In practice, this norm constraint can be enforced via:
\begin{itemize}[leftmargin=1.5em]
  \item Spectral normalization of the block matrix representing \(A\);
  \item Factorizations \(W_{ij} = \alpha_{ij} \tilde W_{ij}\) with
  \(\|\tilde W_{ij}\|_{\mathrm{op}} \le 1\) and
  \(\sum_{j \in \mathcal{N}(i)} |\alpha_{ij}| \le \Lambda_A\) for all \(i\).
\end{itemize}

\subsection{Fairness and Energy Invariants}

As in MOC, we may track per-prime energy and fairness:
\[
  E_i(h) = \|h[i,:]\|_2^2,
  \qquad i=1,\dots,P.
\]
A simple fairness invariant is
\[
  F(h)
  =
  \frac{\min_i E_i(h)}{\max_i E_i(h)},
\]
with the requirement \(F(h) \ge F_{\min}\) for some \(F_{\min} \in (0,1]\). In
the cell, this can be enforced by:
\begin{itemize}[leftmargin=1.5em]
  \item Projecting back into the region \(F(h) \ge F_{\min}\),
  \item Or penalizing violations via an additional loss term.
\end{itemize}
This mirrors the fairness constraint described in MOC's multiplicity space
invariants \cite{MOC}.

\subsection{MOC-Inspired Variants}

The simple definition \eqref{eq:prime-mixing} can be enriched by including
MOC-inspired structure:
\begin{itemize}[leftmargin=1.5em]
  \item \emph{Tier-aware mixing:} define \(\mathcal{N}(i)\) using CRT tiers or
  divisibility patterns (e.g., primes involved in the same \(p^r\)-based tier).
  \item \emph{Permutation-like actions:} in small prime clusters, implement
  \(W_{ij}\) as a convex combination of permutation matrices, paralleling MOC's
  within-cell permutations \(W_p\) \cite{MOC}.
  \item \emph{Attention-based mixing:} implement \(A\) as a local
  self-attention layer over prime channels, with attention restricted to
  \(\mathcal{N}(i)\) and normalized to maintain contractivity.
\end{itemize}

\section{Relationship Between MOC and MultiplicityCell}

\subsection{Carrier Space}

MOC's carrier
\[
  M = (X_n,H)
\]
can be viewed as a finite, discrete realization of
\(H_{\text{lawful}}\), with:
\begin{itemize}[leftmargin=1.5em]
  \item \(X_n\) approximating \(L^2(\mathbb{R})\) on a cyclic lattice;
  \item Hypergraph \(H\) providing a finite relational internal space.
\end{itemize}
The MultiplicityCell's hidden state \(h_t\) compresses this structure into a
finite set of prime channels and internal dimensions.

\subsection{Operators and Blocks}

MOC's prime-indexed operators \((S_p, A_{p^r}, R_{p^r}, W_p, \hat Q_p)\)
naturally map onto the components of \(T_{\Lambda_m}\) in MultiplicityCell:
\begin{itemize}[leftmargin=1.5em]
  \item Subdivision \(S_p\) and rotation \(R_{p^r}\), together with CRT
  projectors \(\Pi_{p^r}\), correspond to the time-sieve block \(B\).
  \item Accents \(A_{p^r}\), permutations \(W_p\), and cross-prime
  noncommutation provide blueprint patterns for the prime-mixing block \(A\).
  \item Relation operators \(\hat Q_p\) and gauge transformations motivate the
  internal block \(E\) and associated projectors.
\end{itemize}

\subsection{Projectors and Certification}

MOC introduces:
\begin{itemize}[leftmargin=1.5em]
  \item Level projectors \(\Pi_{p^r}\),
  \item Spike gates \(\Delta_{p^r}\),
  \item Energy and fairness invariants.
\end{itemize}
These play roles analogous to:
\begin{itemize}[leftmargin=1.5em]
  \item Constitutional projectors \(\Pi_{\text{CSL}}\) (enforcing tier-wise
  structure),
  \item Ethical projector \(P_E\) (enforcing fairness and viability).
\end{itemize}

The resonance functional \(R\) in MOC can be regarded as a specialization of a
more general ``coherence'' functional used to certify lawful trajectories of
the universal recursion \eqref{eq:universal-recursion}.

\section{Example Implementation of the Prime-Mixing Block}

We conclude with a concrete PyTorch-style sketch of the prime-mixing block
\(A\), suitable for embedding within a MultiplicityCell implementation.

\begin{lstlisting}[style=pytorch, caption={Prime-mixing block $A$ in PyTorch-style pseudocode.}]
import torch
import torch.nn as nn
import torch.nn.functional as F

class PrimeMixingBlock(nn.Module):
    def __init__(self, primes, d_int, bandwidth, lambda_A):
        super().__init__()
        self.primes = torch.tensor(primes, dtype=torch.long)
        self.P = len(primes)
        self.d_int = d_int
        self.bandwidth = bandwidth
        self.lambda_A = lambda_A

        # Precompute neighborhoods N(i) based on prime value distance
        # N(i) = { j : |p_j - p_i| <= bandwidth }
        N_list = []
        for i in range(self.P):
            p_i = self.primes[i].item()
            idx = [
                j for j in range(self.P)
                if abs(self.primes[j].item() - p_i) <= bandwidth
            ]
            N_list.append(idx)
        self.N_list = N_list

        # Learnable weight matrices W_ij stored as a single tensor
        # We'll index into it with (i,j)
        self.W = nn.Parameter(
            torch.randn(self.P, self.P, d_int, d_int) * 0.01
        )
        # Bias per prime channel
        self.b = nn.Parameter(torch.zeros(self.P, d_int))

    def forward(self, h):
        """
        h: Tensor of shape (batch_size, P, d_int)
        returns: Tensor of same shape
        """
        B, P, d_int = h.shape
        assert P == self.P and d_int == self.d_int

        # Linear mixing according to neighborhoods
        out = torch.zeros_like(h)
        for i in range(self.P):
            # h_i_new = sum_{j in N(i)} W_ij h_j + b_i
            hij_sum = 0.0
            for j in self.N_list[i]:
                W_ij = self.W[i, j]  # (d_int, d_int)
                hij_sum = hij_sum + torch.einsum(
                    "bd,dk->bk", h[:, j, :], W_ij
                )
            out[:, i, :] = hij_sum + self.b[i]

        # Optional: spectral / Frobenius normalization to enforce ||A|| <= lambda_A.
        # Here we use a simple global rescaling; in practice, use spectral norm.
        with torch.no_grad():
            # Estimate global scale via Frobenius norm of W
            scale = self.W.norm().item()
            if scale > 0:
                factor = min(1.0, self.lambda_A / scale)
            else:
                factor = 1.0
            # Note: for strict control, integrate normalization into parameterization.
        out = out * factor

        return out
\end{lstlisting}

This sketch demonstrates:
\begin{itemize}[leftmargin=1.5em]
  \item How prime neighborhoods can be defined based on arithmetic distance;
  \item How to implement the banded mixing structure \((A h)_i = \sum_{j \in \mathcal{N}(i)} W_{ij} h_j + b_i\);
  \item A naive (but illustrative) way to enforce a global contraction bound
  \(\|A\|_{\mathrm{op}} \le \Lambda_A\). In a production setting, one would
  replace the ad hoc rescaling with a principled spectral normalization.
\end{itemize}

\section{Outlook}

The developments in this thread suggest several natural next steps:
\begin{itemize}[leftmargin=1.5em]
  \item A \emph{constitutional master paper} that defines
  \(H_{\text{lawful}}, U, \Lambda_m, \Pi_{\text{CSL}}, P_E\) in generality and
  then treats MOC, MultiplicityCell, ZetaCell, ethical manifolds, and other
  modules as explicit frames on this common substrate.
  \item A unified \emph{software library} implementing:
  \begin{itemize}
    \item MultiplicityCell (base cell),
    \item ZetaCell (dual-sector cell),
    \item Ethical projectors and invariants,
    \item MOC-style resonance diagnostics and CRT-based tier tools.
  \end{itemize}
  \item A focused \emph{empirical demonstration}, such as prime-locked EEG
  modeling or semantic phase transitions in a small language model, leveraging
  the prime-mixing block \(A\) and resonance-style diagnostics.
\end{itemize}

These steps would further crystallize the unification suggested by MOC and
the broader Multiplicity corpus, and they would provide concrete evidence that
the prime-indexed multiplicity ontology yields practical, falsifiable models.

%====================================================================
% Mathematical Appendix
%====================================================================

\appendix

\section{Mathematical Appendix: Operator Bounds and Fixed-Point Structure}

In this appendix we collect explicit operator norm bounds and fixed-point
results for the prime-mixing block \(A\) and for the one-step update map of
MultiplicityCell, viewed as a finite-dimensional surrogate of the universal
recursion. We assume the notation and constructions of the main text.

\subsection{Notation and Norms}

Let \(P\) denote the number of prime channels and \(d_{\text{int}}\) the
internal dimension per prime. The hidden state of MultiplicityCell at time
\(t\) is
\[
  h_t \in \mathbb{R}^{P \times d_{\text{int}}}.
\]
We will identify \(\mathbb{R}^{P \times d_{\text{int}}}\) with
\(\mathbb{R}^{D}\) where \(D = P \cdot d_{\text{int}}\), either by flattening
along the prime index and internal dimension or via a fixed isometric
identification.

We equip \(\mathbb{R}^{D}\) with the Euclidean norm
\[
  \|h\|_2 = \biggl( \sum_{i=1}^P \sum_{k=1}^{d_{\text{int}}}
  h[i,k]^2 \biggr)^{1/2}.
\]
All operator norms \(\|\cdot\|_{\mathrm{op}}\) in this appendix are with
respect to \(\|\cdot\|_2\).

\subsection{Prime-Mixing Block as a Banded Operator}

Recall the definition of the prime-mixing block \(A\) from the main text:
for \(h \in \mathbb{R}^{P \times d_{\text{int}}}\) and each
\(i \in \{1,\dots,P\}\),
\begin{equation}
  (A h)_i
  =
  \sum_{j \in \mathcal{N}(i)} W_{ij} \, h[j,:] + b_i,
  \label{eq:appendix-Ah}
\end{equation}
where:
\begin{itemize}
  \item \(\mathcal{N}(i) \subseteq \{1,\dots,P\}\) is the neighborhood of
  prime index \(i\);
  \item \(W_{ij} \in \mathbb{R}^{d_{\text{int}}\times d_{\text{int}}}\) are
  learnable weight matrices, with \(W_{ij} = 0\) for \(j \notin \mathcal{N}(i)\);
  \item \(b_i \in \mathbb{R}^{d_{\text{int}}}\) is a bias vector.
\end{itemize}

We first bound \(\|A\|_{\mathrm{op}}\) in terms of local norms of the blocks
\(W_{ij}\) and the neighborhood sizes.

\begin{lemma}[Row-wise block bound]
\label{lem:row-wise-block-bound}
Let \(h \in \mathbb{R}^{P \times d_{\text{int}}}\) and \(A\) be defined by
\eqref{eq:appendix-Ah} with \(b_i = 0\) for simplicity.
For each \(i\), set
\[
  \alpha_i
  =
  \sum_{j \in \mathcal{N}(i)} \|W_{ij}\|_{\mathrm{op}},
  \qquad
  \alpha_{\max} = \max_{1 \le i \le P} \alpha_i.
\]
Then for all \(h\),
\[
  \|A h\|_2 \le \alpha_{\max} \|h\|_2 + \beta \|h\|_2,
\]
where \(\beta \ge 0\) is an explicit term that depends only on overlaps between
neighborhoods \(\mathcal{N}(i)\). In particular, if the neighborhoods are
uniformly bounded and have limited overlap, one has
\[
  \|A\|_{\mathrm{op}} \le C \alpha_{\max}
\]
for some constant \(C\) depending only on the neighborhood structure.
\end{lemma}

\begin{proof}
We write the hidden state as \(h = (h[1,:],\dots,h[P,:])\), with each
\(h[i,:]\in\mathbb{R}^{d_{\text{int}}}\). By \eqref{eq:appendix-Ah} with
\(b_i=0\),
\[
  (A h)_i = \sum_{j\in\mathcal{N}(i)} W_{ij} h[j,:].
\]
By the triangle inequality and submultiplicativity of \(\|\cdot\|_2\),
\[
  \|(A h)_i\|_2
  \le
  \sum_{j\in\mathcal{N}(i)} \|W_{ij}\|_{\mathrm{op}} \, \|h[j,:]\|_2
  =
  \sum_{j\in\mathcal{N}(i)} a_{ij} \, \|h[j,:]\|_2
\]
where \(a_{ij} = \|W_{ij}\|_{\mathrm{op}}\).

Let \(u_i = \|(A h)_i\|_2\) and \(v_j = \|h[j,:]\|_2\).
We thus obtain
\[
  u_i \le \sum_{j\in\mathcal{N}(i)} a_{ij} v_j.
\]
Squaring and summing over \(i\), we have
\[
  \|A h\|_2^2
  =
  \sum_{i=1}^P u_i^2
  \le
  \sum_{i=1}^P
  \biggl( \sum_{j\in\mathcal{N}(i)} a_{ij} v_j \biggr)^2.
\]
A straightforward but slightly tedious estimate using Cauchy--Schwarz and the
boundedness of \(|\mathcal{N}(i)|\) and neighborhood overlaps yields a bound
of the form
\[
  \sum_{i=1}^P
  \biggl( \sum_{j\in\mathcal{N}(i)} a_{ij} v_j \biggr)^2
  \le
  C^2 \alpha_{\max}^2 \sum_{j=1}^P v_j^2,
\]
where \(C\) is a constant depending only on the structure of the neighborhoods
\(\mathcal{N}(i)\). Taking square roots yields
\[
  \|A h\|_2 \le C \alpha_{\max} \|h\|_2.
\]
The displayed inequality with the additional term \(\beta \|h\|_2\) is a more
general statement where contributions from overlaps can be explicitly grouped,
but for our purposes it suffices to observe the existence of a constant
\(C \ge 1\) bounding the operator norm as claimed.
\end{proof}

\begin{corollary}[Contractive parametric regime]
\label{cor:contractive-A}
Suppose the neighborhoods \(\{\mathcal{N}(i)\}\) are fixed and satisfy the
hypotheses of Lemma~\ref{lem:row-wise-block-bound} with constant \(C\). If we
impose the constraint
\[
  \alpha_{\max}
  =
  \max_i \sum_{j\in\mathcal{N}(i)} \|W_{ij}\|_{\mathrm{op}}
  \le \frac{\Lambda_A}{C}
\]
for some \(\Lambda_A \in (0,1)\), then
\[
  \|A\|_{\mathrm{op}} \le \Lambda_A.
\]
\end{corollary}

\begin{proof}
Immediate from Lemma~\ref{lem:row-wise-block-bound} and the choice of
\(\alpha_{\max}\).
\end{proof}

Thus, by controlling the block norms \(\|W_{ij}\|_{\mathrm{op}}\) and the
neighborhood structure, we obtain an explicit contractivity bound for \(A\).

\subsection{Bias Terms and Affine Maps}

In the presence of bias terms \(b_i\), we have an affine map
\[
  A h + b,
  \qquad
  b = (b_1,\dots,b_P).
\]
The operator norm bound in Corollary~\ref{cor:contractive-A} remains valid for
the linear part of the map; bias terms do not affect the Lipschitz constant
with respect to \(\|\cdot\|_2\). In particular, the mapping
\[
  h \mapsto A h + b
\]
is \(\Lambda_A\)-Lipschitz if \(\|A\|_{\mathrm{op}} \le \Lambda_A\).

\subsection{One-Step Cell Map and Fixed Points}

We now consider the full one-step update of MultiplicityCell, in a simplified
setting where the input \(x_t\) is either fixed or does not affect the
Lipschitz constant. Let \(T_{\Lambda_m}\) be decomposed as
\[
  T_{\Lambda_m}
  =
  A + B + E,
\]
where \(A\) is the prime-mixing block, \(B\) is a time-sieve block, and \(E\)
an internal block. We assume that each of these blocks is Lipschitz with
constants \(\Lambda_A\), \(\Lambda_B\), \(\Lambda_E\), respectively:
\[
  \|A h - A h'\|_2 \le \Lambda_A \|h - h'\|_2, \quad
  \|B h - B h'\|_2 \le \Lambda_B \|h - h'\|_2, \quad
  \|E h - E h'\|_2 \le \Lambda_E \|h - h'\|_2.
\]
We also assume \(\Lambda_A,\Lambda_B,\Lambda_E \ge 0\).

\begin{lemma}[Lipschitz bound for $T_{\Lambda_m}$]
\label{lem:Lipschitz-T}
Under the above assumptions,
\[
  \|T_{\Lambda_m}(h) - T_{\Lambda_m}(h')\|_2
  \le
  (\Lambda_A + \Lambda_B + \Lambda_E)\, \|h - h'\|_2.
\]
\end{lemma}

\begin{proof}
By linearity (or at least additivity) of \(A\), \(B\), \(E\),
\[
  T_{\Lambda_m}(h) - T_{\Lambda_m}(h')
  =
  (A h - A h') + (B h - B h') + (E h - E h').
\]
By the triangle inequality and the assumed Lipschitz bounds,
\[
  \|T_{\Lambda_m}(h) - T_{\Lambda_m}(h')\|_2
  \le
  \|A h - A h'\|_2 + \|B h - B h'\|_2 + \|E h - E h'\|_2
  \le
  (\Lambda_A + \Lambda_B + \Lambda_E) \|h - h'\|_2.
\]
\end{proof}

Next we incorporate the projectors \(\Pi_{\text{CSL}}\) and \(P_E\). For
simplicity we assume they are non-expansive maps (i.e., 1-Lipschitz) with
respect to \(\|\cdot\|_2\).

\begin{assumption}[Non-expansive projectors]
\label{ass:nonexpansive-projectors}
The projectors \(\Pi_{\text{CSL}}\) and \(P_E\) satisfy
\[
  \|\Pi_{\text{CSL}}(h) - \Pi_{\text{CSL}}(h')\|_2 \le \|h - h'\|_2,
  \qquad
  \|P_E(h) - P_E(h')\|_2 \le \|h - h'\|_2
\]
for all \(h,h'\).
\end{assumption}

Define the full one-step cell map
\[
  F(h)
  =
  P_E\bigl(
    \Pi_{\text{CSL}} T_{\Lambda_m}(h)
  \bigr).
\]

\begin{theorem}[Global Lipschitz bound and contraction regime]
\label{thm:cell-contraction}
Under Assumption~\ref{ass:nonexpansive-projectors} and the Lipschitz
assumptions above, the map \(F\) satisfies
\[
  \|F(h) - F(h')\|_2
  \le
  L_F \|h - h'\|_2,
  \qquad
  L_F = \Lambda_A + \Lambda_B + \Lambda_E.
\]
In particular, if
\[
  \Lambda_A + \Lambda_B + \Lambda_E < 1,
\]
then \(F\) is a contraction on \(\mathbb{R}^{D}\) and admits a unique fixed
point \(h^\star\) for any fixed input (or in the autonomous case).
\end{theorem}

\begin{proof}
Let \(h,h'\in\mathbb{R}^D\). By Assumption~\ref{ass:nonexpansive-projectors},
\[
  \|F(h) - F(h')\|_2
  =
  \|P_E(\Pi_{\text{CSL}} T_{\Lambda_m}(h))
    - P_E(\Pi_{\text{CSL}} T_{\Lambda_m}(h'))\|_2
  \le
  \|\Pi_{\text{CSL}} T_{\Lambda_m}(h)
    - \Pi_{\text{CSL}} T_{\Lambda_m}(h')\|_2
  \le
  \|T_{\Lambda_m}(h) - T_{\Lambda_m}(h')\|_2.
\]
Applying Lemma~\ref{lem:Lipschitz-T} yields
\[
  \|F(h) - F(h')\|_2
  \le
  (\Lambda_A + \Lambda_B + \Lambda_E) \|h - h'\|_2
  = L_F \|h - h'\|_2.
\]
If \(L_F < 1\), then \(F\) is a contraction on the complete metric space
\((\mathbb{R}^D,\|\cdot\|_2)\). By the Banach fixed-point theorem, \(F\) has a
unique fixed point \(h^\star\), and iterates \(h_{t+1} = F(h_t)\) converge to
\(h^\star\) from any initial condition.
\end{proof}

In practice, the universal multiplicity constant \(\Lambda_m\) can be chosen
so that
\[
  \Lambda_A + \Lambda_B + \Lambda_E \le \Lambda_m < 1,
\]
with each block's parameters constrained accordingly (e.g. via
Corollary~\ref{cor:contractive-A} for \(A\)).

\subsection{Fairness Invariant Under Small Mixing}

MOC introduces a fairness invariant
\[
  F(h) = \frac{\min_i E_i(h)}{\max_i E_i(h)},
  \qquad
  E_i(h) = \|h[i,:]\|_2^2,
\]
to quantify balance among channels \cite{MOC}. Here we show that, under a
small mixing regime, the fairness ratio is approximately preserved by the
prime-mixing block \(A\).

\begin{lemma}[Perturbative preservation of fairness]
\label{lem:fairness-preservation}
Let \(h \in \mathbb{R}^{P \times d_{\text{int}}}\) and suppose
\(E_i(h) = \|h[i,:]\|_2^2 > 0\) for all \(i\). Let \(A\) be a prime-mixing
block satisfying
\[
  \|A\|_{\mathrm{op}} \le \varepsilon
\]
for some \(\varepsilon \in (0,1)\). Consider the updated state
\[
  \tilde h = h + A h.
\]
Then there exists a constant \(C\) (depending on \(P\) and
\(d_{\text{int}}\)) such that
\[
  |F(\tilde h) - F(h)|
  \le
  C \varepsilon.
\]
In particular, for sufficiently small \(\varepsilon\), the fairness ratio
\(F(\tilde h)\) remains within a controlled perturbation of \(F(h)\).
\end{lemma}

\begin{proof}
Let \(\delta h = A h\). Then \(\|\delta h\|_2 \le \varepsilon \|h\|_2\) by
assumption. For a fixed index \(i\),
\[
  \|\tilde h[i,:]\|_2^2
  =
  \|h[i,:] + \delta h[i,:]\|_2^2
  =
  \|h[i,:]\|_2^2
  + 2 \langle h[i,:], \delta h[i,:]\rangle
  + \|\delta h[i,:]\|_2^2.
\]
By Cauchy--Schwarz,
\[
  |\langle h[i,:], \delta h[i,:]\rangle|
  \le
  \|h[i,:]\|_2 \|\delta h[i,:]\|_2
  \le
  \|h[i,:]\|_2 \|\delta h\|_2
  \le
  \|h[i,:]\|_2 \varepsilon \|h\|_2.
\]
Similarly,
\[
  \|\delta h[i,:]\|_2^2
  \le
  \|\delta h\|_2^2
  \le
  \varepsilon^2 \|h\|_2^2.
\]
Therefore,
\[
  E_i(\tilde h)
  =
  \|\tilde h[i,:]\|_2^2
  =
  E_i(h)
  + \mathcal{O}(\varepsilon \|h[i,:]\|_2 \|h\|_2)
  + \mathcal{O}(\varepsilon^2 \|h\|_2^2).
\]
The constants implicit in \(\mathcal{O}(\cdot)\) depend only on dimension.

Consider now the minimum and maximum energies. Let
\[
  E_{\min} = \min_i E_i(h), \quad
  E_{\max} = \max_i E_i(h),
  \quad
  F(h) = \frac{E_{\min}}{E_{\max}}.
\]
For \(\varepsilon\) sufficiently small, perturbation bounds on the minima and
maxima of a finite family of real numbers yield
\[
  E_{\min}(\tilde h)
  =
  E_{\min} + \mathcal{O}(\varepsilon \|h\|_2^2), \quad
  E_{\max}(\tilde h)
  =
  E_{\max} + \mathcal{O}(\varepsilon \|h\|_2^2),
\]
with constants depending on \(P\) and \(d_{\text{int}}\). Consequently,
\[
  F(\tilde h)
  =
  \frac{E_{\min}(\tilde h)}{E_{\max}(\tilde h)}
  =
  \frac{E_{\min}}{E_{\max}}
  + \mathcal{O}(\varepsilon)
  =
  F(h) + \mathcal{O}(\varepsilon).
\]
This yields the claimed bound with some \(C>0\).
\end{proof}

Lemma~\ref{lem:fairness-preservation} shows that, if the prime-mixing block
\(A\) is kept in a sufficiently contractive regime (small \(\varepsilon\)), the
fairness invariant remains approximately stable. In the context of the global
framework, the ethical projector \(P_E\) can be designed to correct any
residual deviations from a target fairness threshold, yielding an overall
fairness-preserving dynamics.

\subsection{Remarks on MOC Resonance and Operator Bounds}

While this appendix focuses on MultiplicityCell, analogous operator norm
considerations apply to MOC's operator words. Given a composite word
\[
  W = \hat O_m \cdots \hat O_1
\]
with each \(\hat O_k\) acting on \(X_n\) (and/or \(H\)), one can bound
\(\|W\|_{\mathrm{op}}\) as
\[
  \|W\|_{\mathrm{op}}
  \le
  \prod_{k=1}^m \|\hat O_k\|_{\mathrm{op}},
\]
and choose operator families and weights such that this product remains
controlled. The resonance functional \(R\) can then be related to these bounds
in specific settings, e.g., when certain operators are unitary (rotations,
permutations) and others are small perturbations (accents, gates) \cite{MOC}.

Developing a full theory of operator norm bounds for MOC's joint action on
\((X_n,H)\) and their relation to resonance is a natural direction for further
appendix-level analysis.


\nocite{*}
\bibliographystyle{unsrt}  
\bibliography{references}
\end{document}